\chapter{LITERATURE SURVEY}
\begingroup
\hyphenpenalty=10000
\exhyphenpenalty=10000
\sloppy

This chapter presents an extensive review of existing research in the domain of autoscaling in cloud computing and microservices-based systems, with a particular focus on performance-aware and queuing-theory-based approaches. Traditional autoscaling techniques, which rely on static CPU and memory utilization thresholds, are increasingly inadequate in modern cloud environments due to highly dynamic workloads and complex service interactions \cite{khaleq2021, fourati2021}.

Recent advancements in performance modeling and intelligent scaling mechanisms have introduced more adaptive solutions that consider workload characteristics such as request arrival rate, service time, and system utilization \cite{wang2024}. Queuing theory-based models in particular, provide a more accurate representation of system behavior by capturing the relationship between demand and service capacity, enabling more informed scaling decisions \cite{lai2025}.

The reviewed literature is systematically organized by categorizing existing works into threshold-based autoscaling methods, predictive and machine learning-based approaches and queuing-theory-driven scaling techniques. Each category is analyzed to highlight its strengths, limitations, and research gaps, which form the foundation and motivation for the proposed intelligent autoscaling system.\\
\\
\section{EXISTING WORK}

Existing autoscaling research primarily focuses on three approaches: threshold-based methods that react to CPU/memory utilization but suffer from delayed responses \cite{fourati2021}, predictive machine learning approaches that forecast demands but require extensive training data \cite{wang2024}, and queuing theory-based models that provide mathematical foundations for performance prediction but lack practical implementation in real microservices environments \cite{lai2025}. Current solutions fail to integrate theoretical modeling with real-time decision-making, optimize scaling overhead, or coordinate scaling across interdependent services, revealing the need for an intelligent framework that bridges theory and practice \cite{khaleq2021}.

\subsection{Threshold-Based Autoscaling Approaches}

Early autoscaling systems predominantly relied on threshold-based mechanisms, where scaling decisions are triggered based on predefined CPU and memory utilization limits \cite{khaleq2021}. These systems operate on simple conditional logic, initiating scale-out operations when resource utilization exceeds upper thresholds (typically 70-80%) and scale-in operations when utilization falls below lower thresholds (typically 30-50%). While these approaches are simple and computationally efficient, they fail to capture the actual system performance under dynamic workloads, as resource utilization alone does not accurately reflect service quality or user experience.

Several cloud platforms implement reactive autoscaling policies that scale resources when utilization crosses fixed thresholds \cite{fourati2021}. Major cloud providers such as AWS Auto Scaling, Google Cloud Autoscaler, and Azure Virtual Machine Scale Sets employ variations of this approach, often combined with cooldown periods to prevent oscillatory scaling behavior. Although effective under stable conditions with predictable traffic patterns, these methods struggle with bursty workloads and often result in delayed scaling actions leading to increased response time and temporary system overload.

The fundamental limitation of threshold-based approaches lies in their reactive nature and oversimplified performance indicators. Additionally, they do not consider request-level metrics such as arrival rate and service time, limiting their ability to model real system behavior. This becomes particularly problematic in microservices architectures where service interdependencies can create bottlenecks that are not reflected in individual service resource utilization metrics. Furthermore, the static nature of thresholds means these systems cannot adapt to changing workload characteristics or application behavior over time, resulting in suboptimal scaling decisions that either waste resources through over-provisioning or compromise performance through under-provisioning.

\subsection{Predictive and Machine Learning-Based Autoscaling}

To overcome the limitations of static threshold-based methods, recent research has explored predictive and machine learning-based autoscaling techniques for microservices architectures \cite{wang2024}. These approaches leverage historical workload data from distributed systems to forecast future demand patterns and make proactive scaling decisions, particularly relevant for applications with cyclical traffic such as course registration systems where enrollment periods create predictable load spikes.

Machine learning models such as regression analysis for capacity planning, time-series forecasting techniques like LSTM networks for predicting registration traffic patterns, and reinforcement learning algorithms for adaptive resource allocation have been applied to predict traffic patterns and optimize resource allocation across microservices chains. Some implementations focus on learning the interdependencies between services in multi-tier architectures, where authentication service load can predict downstream demand for course catalog and seat reservation services.

While these approaches demonstrate improved responsiveness compared to reactive methods in handling bursty academic workloads, they face significant challenges in microservices environments. They require extensive training data spanning multiple registration cycles and diverse traffic scenarios, which may not capture the unique characteristics of different academic institutions or course types. The models may not generalize well under highly unpredictable workloads such as emergency course changes, system migrations, or unprecedented enrollment surges that fall outside historical patterns.

\subsection{Queuing Theory-Based Autoscaling Models}

Queuing theory provides a more analytical approach to autoscaling by modeling each service as a queuing system with mathematically defined characteristics. These models consider key parameters such as arrival rate (λ), service rate (μ), and number of servers (c) to estimate performance metrics like utilization (ρ), waiting time, and response time \cite{lai2025}. This theoretical foundation enables precise prediction of system behavior under various load conditions, making it particularly suitable for microservices architectures where service interdependencies create complex performance relationships.

Several studies have demonstrated that queuing models, particularly M/M/c systems representing multi-server configurations, can provide more accurate insights into system behavior compared to traditional threshold-based methods \cite{khaleq2021}. The M/M/c model assumes Poisson arrival processes and exponential service times, which closely approximate real-world web service behavior in many scenarios. Research has shown that these models can accurately predict queue lengths, response times, and optimal server configurations for maintaining desired service levels \cite{fourati2021}.\\

Advanced queuing theory applications have explored multi-class queuing networks to model complex microservices interactions, where different request types (such as authentication, course lookup, and seat reservation) have varying service requirements and priorities \cite{wang2024}. Some implementations have incorporated queuing network models to analyze end-to-end performance across service chains, providing insights into bottleneck identification and capacity planning.

However, many existing implementations use these models in isolation for performance analysis and capacity planning, but do not integrate them into real-time decision-making systems for dynamic autoscaling. The gap between theoretical modeling and practical implementation remains significant, as most queuing theory applications focus on offline analysis rather than online control systems. As a result, their practical applicability in dynamic cloud environments remains limited, despite their superior accuracy in modeling system behavior compared to simple threshold-based approaches. The challenge lies in translating mathematical predictions into actionable scaling decisions while maintaining the computational efficiency required for real-time operations.

\subsection{Adaptive and Multi-Factor Autoscaling Frameworks}

Recent research has focused on combining multiple metrics and decision strategies to improve autoscaling performance in complex microservices environments \cite{arulmary2026}. Multi-factor approaches consider a combination of system-level metrics (CPU, memory, network I/O) and application-level metrics (request latency, throughput, error rates) to make more informed scaling decisions. These hybrid frameworks recognize that single-metric approaches are insufficient for capturing the multidimensional nature of microservices performance, particularly in scenarios like course registration systems where user experience depends on end-to-end response times across multiple service interactions.

Some implementations have explored weighted scoring systems that combine different performance indicators, adaptive threshold mechanisms that adjust based on historical patterns, and rule-based engines that encode domain-specific knowledge about application behavior \cite{fourati2021}. Advanced hybrid approaches have incorporated feedback control loops that continuously adjust scaling parameters based on observed system performance, attempting to create self-tuning autoscaling systems that adapt to changing workload characteristics.

Container orchestration platforms like Kubernetes have introduced Horizontal Pod Autoscalers (HPA) and Vertical Pod Autoscalers (VPA) that support custom metrics beyond basic resource utilization, enabling more sophisticated scaling policies \cite{wang2024}. These systems allow integration of application-specific metrics such as queue depth, active connections, or business-level indicators like concurrent user sessions in registration systems.

While these frameworks demonstrate improved decision accuracy compared to single-metric approaches, they often lack a structured mathematical foundation and may rely on heuristic rules that are difficult to optimize or validate. The combination of multiple metrics frequently involves ad-hoc weighting schemes without theoretical justification, making it challenging to predict system behavior or guarantee performance objectives. Additionally, these approaches do not effectively address optimal replica allocation or dynamic traffic distribution across services, which are critical for maintaining system stability in microservices architectures where service interdependencies can create cascading performance issues. The absence of coordinated scaling decisions across the entire service chain often results in resource imbalances that can negate the benefits of sophisticated single-service scaling policies.

\section{OBSERVATIONS FROM THE EXISTING WORK}

The reviewed literature indicates that traditional threshold-based autoscaling methods are simple to implement and computationally efficient but insufficient for handling dynamic and bursty workloads characteristic of modern cloud applications \cite{khaleq2021}. These reactive approaches suffer from inherent delays in detection and response, often resulting in performance degradation during traffic spikes and inefficient resource utilization during normal operations. The static nature of threshold values also prevents adaptation to changing application behavior and workload patterns over time.

Predictive and machine learning-based approaches demonstrate improved responsiveness by anticipating future resource demands, but often require large amounts of training data spanning multiple workload cycles and lack interpretability in their decision-making processes \cite{wang2024}. These methods face challenges in generalizing to unprecedented traffic patterns and may introduce computational overhead that delays scaling decisions. Furthermore, the black-box nature of many machine learning models makes it difficult to understand and validate scaling decisions, particularly in critical applications where performance guarantees are essential.

Queuing theory-based models provide a more accurate and mathematically rigorous representation of system behavior by capturing the fundamental relationship between arrival rate, service rate, and system performance metrics such as utilization and response time \cite{lai2025}. These analytical approaches offer the advantage of theoretical soundness and predictable behavior, enabling precise performance predictions under various load conditions. However, their practical implementation in real-time autoscaling systems remains limited due to the complexity of translating mathematical models into operational control mechanisms.

However, many existing approaches either rely on reactive threshold mechanisms that respond only after performance degradation occurs, or use predictive models without integrating real-time analytical decision-making capabilities that can provide immediate insights into system state. Several frameworks do not fully utilize the potential of queuing models for continuous scaling decisions and often ignore key performance metrics such as response time and queue length that are critical for maintaining service quality. Additionally, most systems focus primarily on scaling individual service resources in isolation and do not address optimal replica allocation strategies or efficient traffic distribution mechanisms across interdependent microservices \cite{arulmary2026}. This limitation becomes particularly problematic in service chain architectures where bottlenecks in one component can cascade throughout the entire system, requiring coordinated scaling decisions that consider the holistic system performance rather than individual service metrics.

\section{LIMITATIONS OF EXISTING WORK}

Existing autoscaling methods are largely reactive in nature and fail to handle dynamic workloads effectively, particularly in microservices environments where service interdependencies amplify performance issues \cite{khaleq2021}. These threshold-based approaches suffer from inherent delays in detection and response, often resulting in performance degradation during traffic spikes and suboptimal resource utilization during normal operations.

Predictive approaches show promise in anticipating resource demands but depend heavily on historical data and lack adaptability to unprecedented traffic patterns or system changes \cite{wang2024}. The computational overhead of complex prediction models can also introduce latency into scaling decisions, while their black-box nature makes it difficult to validate and optimize scaling policies for specific application requirements.

Many existing systems do not fully utilize the mathematical rigor of queuing theory for real-time decision-making or fail to address critical aspects such as optimal replica allocation and coordinated traffic distribution across microservices. This leads to inefficient resource usage, unbalanced load distribution, and reduced overall system performance, particularly in service chain architectures where bottlenecks can cascade throughout the entire system. The lack of holistic approaches that consider both individual service performance and system-wide optimization represents a significant gap in current autoscaling research.

\section{CHALLENGES}

Key challenges identified include handling highly dynamic and bursty workloads characteristic of applications like course registration systems, where traffic can spike dramatically within minutes during enrollment periods. Making accurate real-time scaling decisions that balance performance requirements with resource costs presents another significant challenge, particularly when traditional metrics fail to capture actual system behavior under varying load conditions.

Avoiding delayed or inefficient scaling responses is critical, as the time required to detect performance issues, make scaling decisions, and provision new resources often exceeds the duration of traffic spikes, rendering reactive approaches ineffective. The challenge is compounded by the need to prevent oscillatory scaling behavior that can destabilize system performance.

Ensuring optimal resource utilization while maintaining strict SLA requirements creates a complex optimization problem, where under-provisioning leads to performance degradation and over-provisioning results in unnecessary costs. This balance becomes more difficult in microservices architectures where different services have varying resource requirements and performance characteristics.

Additionally, achieving efficient load distribution across interdependent microservices without introducing significant computational overhead remains a significant challenge. The need to coordinate scaling decisions across service chains while maintaining real-time responsiveness requires sophisticated algorithms that can process multiple performance metrics simultaneously. Furthermore, ensuring that scaling actions in one service do not create bottlenecks in dependent services adds another layer of complexity to the autoscaling problem.

\section{SUMMARY OF LITERATURE SURVEY}

This chapter reviewed existing autoscaling approaches with a focus on threshold-based, predictive, and queuing-theory-driven methods. While prior works improve scalability and responsiveness, they exhibit limitations in handling dynamic workloads, real-time decision-making, and efficient resource utilization.

These gaps motivate the proposed work, which introduces a queuing-theory-based intelligent autoscaling system using Multistage Fine-Grained Dynamic Scaling (MFDS), Microservice Scaling Optimization (MSO), and Fair Probabilistic Routing (FPR) to achieve efficient scaling, balanced load distribution, and improved system performance under dynamic conditions[1].